% GENERATED FILE. Edit canonical lessons, then run npm run generate:books.
% canonical-source-hash: fnv1a64:e73eb59ea01a5d75
% canonical-lessons: LA-C121-cochlear, LA-C121-olla, LA-C121-patina, LA-C121-catinus, LA-C121-candela

\chapter{Spoon, Cooking pot, Dish, Bowl, Candle}
\label{ch:la-spoon-cooking-pot-dish-bowl-candle}
\hlchaptermodality{121}

\begin{chapteropening}
\textbf{By the end of this chapter:} \emph{I can say a spoon, a cooking pot, a dish, a bowl and a candle.}

Complete the last lesson of chapter 121: I can say a spoon, a cooking pot, a dish, a bowl and a candle.
\end{chapteropening}

\section[cochlear, cochleāris]{cochlear, cochleāris — a spoon}

\label{lesson:LA-C121-cochlear}



\begin{quote}
Before the new one: say the Latin for a blanket, then the Latin for a napkin.
\end{quote}

\subsection*{You'll want to know: cochlear, cochleāris}
\textbf{cochlear, cochleāris} — \textquotedblleft{}a spoon\textquotedblright{}.

\begin{cousinweb}
\textbf{cochlear} is built on \emph{cochlea}, \textquotedblleft{}a snail\textquotedblright{} (EtymDB 2.1), a word Latin took from Greek \emph{kokhlias}, \textquotedblleft{}snail, spiral\textquotedblright{}: the spoon was first for eating snails, its pointed end for picking them out.
\end{cousinweb}

\begin{grammarlens}[title={What you've built}]
One more word for this chapter.
\end{grammarlens}

\subsection*{Your turn}
\begin{itemize}
\raggedright
  \item \emph{Say it:} \emph{cochlear}
  \item \emph{Say it:} \emph{cochlear}, once more
  \item \emph{Say it:} say \emph{cochlear}
  \item \emph{Recall:} say the Latin for a blanket, then the Latin for a napkin, then say \emph{cochlear} again
\end{itemize}

\begin{tcolorbox}[breakable,colback=teal!4,colframe=teal!35!black,title={Before you move on}]
What does cochlear, cochleāris mean? (\textquotedblleft{}A spoon\textquotedblright{}.) Say it once more.
\end{tcolorbox}
\section[ōlla, ōllae]{ōlla, ōllae — a cooking pot}

\label{lesson:LA-C121-olla}



\begin{quote}
Before the new one: say the Latin for a napkin, then the Latin for a spoon.
\end{quote}

\subsection*{You'll want to know: ōlla, ōllae}
\textbf{ōlla, ōllae} — \textquotedblleft{}a cooking pot\textquotedblright{}.

\begin{cousinweb}
\textbf{ōlla} is inherited, usually traced through Proto-Italic \emph{*aukslā}, \textquotedblleft{}pot\textquotedblright{}, to Proto-Indo-European \emph{*Hukwsleh2}, \textquotedblleft{}cooking pot\textquotedblright{} (EtymDB 2.1). An older Latin spelling was \emph{aula} (Lewis \& Short).
\end{cousinweb}

\begin{grammarlens}[title={What you've built}]
One more word for this chapter.
\end{grammarlens}

\subsection*{Your turn}
\begin{itemize}
\raggedright
  \item \emph{Say it:} \emph{ōlla}
  \item \emph{Say it:} \emph{ōlla}, once more
  \item \emph{Say it:} say \emph{ōlla}
  \item \emph{Recall:} say the Latin for a napkin, then the Latin for a spoon, then say \emph{ōlla} again
\end{itemize}

\begin{tcolorbox}[breakable,colback=teal!4,colframe=teal!35!black,title={Before you move on}]
What does ōlla, ōllae mean? (\textquotedblleft{}A cooking pot\textquotedblright{}.) Say it once more.
\end{tcolorbox}
\section[patina, patinae]{patina, patinae — a dish, a pan}

\label{lesson:LA-C121-patina}



\begin{quote}
Before the new one: say the Latin for a spoon, then the Latin for a cooking pot.
\end{quote}

\subsection*{You'll want to know: patina, patinae}
\textbf{patina, patinae} — \textquotedblleft{}a dish, a pan\textquotedblright{}.

\begin{cousinweb}
\textbf{patina} was probably borrowed from Greek \emph{patanē}, \textquotedblleft{}a flat dish\textquotedblright{} (EtymDB 2.1); Lewis \& Short's link with \emph{pateō}, \textquotedblleft{}to lie open\textquotedblright{}, is less likely. English \emph{patina} is the same word.
\end{cousinweb}

\begin{grammarlens}[title={What you've built}]
One more word for this chapter.
\end{grammarlens}

\subsection*{Your turn}
\begin{itemize}
\raggedright
  \item \emph{Say it:} \emph{patina}
  \item \emph{Say it:} \emph{patina}, once more
  \item \emph{Say it:} say \emph{patina}
  \item \emph{Recall:} say the Latin for a spoon, then the Latin for a cooking pot, then say \emph{patina} again
\end{itemize}

\begin{tcolorbox}[breakable,colback=teal!4,colframe=teal!35!black,title={Before you move on}]
What does patina, patinae mean? (\textquotedblleft{}A dish, a pan\textquotedblright{}.) Say it once more.
\end{tcolorbox}
\section[catīnus, catīnī]{catīnus, catīnī — a bowl, a deep dish}

\label{lesson:LA-C121-catinus}



\begin{quote}
Before the new one: say the Latin for a cooking pot, then the Latin for a dish.
\end{quote}

\subsection*{You'll want to know: catīnus, catīnī}
\textbf{catīnus, catīnī} — \textquotedblleft{}a bowl, a deep dish\textquotedblright{}.

\begin{cousinweb}
\textbf{catīnus} has no sure origin; Varro compared a Sicilian Greek \emph{katinon}, \textquotedblleft{}dish\textquotedblright{} (Lewis \& Short), and it may be a loan from there.
\end{cousinweb}

\begin{grammarlens}[title={What you've built}]
One more word for this chapter.
\end{grammarlens}

\subsection*{Your turn}
\begin{itemize}
\raggedright
  \item \emph{Say it:} \emph{catīnus}
  \item \emph{Say it:} \emph{catīnus}, once more
  \item \emph{Say it:} say \emph{catīnus}
  \item \emph{Recall:} say the Latin for a cooking pot, then the Latin for a dish, then say \emph{catīnus} again
\end{itemize}

\begin{tcolorbox}[breakable,colback=teal!4,colframe=teal!35!black,title={Before you move on}]
What does catīnus, catīnī mean? (\textquotedblleft{}A bowl, a deep dish\textquotedblright{}.) Say it once more.
\end{tcolorbox}
\section[candēla, candēlae]{candēla, candēlae — a candle}

\label{lesson:LA-C121-candela}



\begin{quote}
Before the new one: say the Latin for a dish, then the Latin for a bowl.

Put three hours in order, first to last: \emph{prīma hōra}, \emph{sexta hōra}, \emph{nōna hōra}. Which of them is \emph{ante merīdiem}?
\end{quote}

\subsection*{You'll want to know: candēla, candēlae}
\textbf{candēla, candēlae} — \textquotedblleft{}a candle\textquotedblright{}.

\begin{cousinweb}
\textbf{candēla} is built on \emph{candeō}, \textquotedblleft{}to shine, to glow white\textquotedblright{} (Lewis \& Short), with \emph{-ēla}, an ending for a tool. \emph{candeō} is usually traced to Proto-Indo-European \emph{*(s)kand-}, \textquotedblleft{}to shine\textquotedblright{} (EtymDB 2.1). English \emph{candle} comes from \emph{candēla}.
\end{cousinweb}

\begin{grammarlens}[title={What you've built}]
One more word for this chapter.
\end{grammarlens}

\subsection*{Your turn}
\begin{itemize}
\raggedright
  \item \emph{Say it:} \emph{candēla}
  \item \emph{Say it:} \emph{candēla}, once more
  \item \emph{Say it:} say \emph{candēla}
  \item \emph{Recall:} say the Latin for a dish, then the Latin for a bowl, then say \emph{candēla} again
\end{itemize}

\begin{tcolorbox}[breakable,colback=teal!4,colframe=teal!35!black,title={Before you move on}]
What does candēla, candēlae mean? (\textquotedblleft{}A candle\textquotedblright{}.) Say it once more.
\end{tcolorbox}
